\documentclass[11pt]{article}
\usepackage{preamble}
\usepackage{subfiles} % load last: enables per-section standalone compilation

% Previous working title: "Braid Arrangements for Concurrent Programs".
\title{Complexified Cube Chains and a Braid Invariant of Precubical Sets}
\author{% TODO: author name(s) and affiliation.  Left blank deliberately ---
        % do not circulate until this is filled in.
  }
\date{\today}

\begin{document}
\maketitle

\begin{abstract}
Ziemiański's cube chain category $\Ch(K)$ computes the space of directed paths on a precubical
set $K$: its nerve is naturally weakly homotopy equivalent to $\dP(K)_0^1$
\cite[Thm.~7.6]{Z2}. Paliga and Ziemiański used this to define braid-group invariants of
precubical sets, and observed that those invariants vanish on the class of greatest interest to
concurrency --- sculptures, equivalently Euclidean cubical complexes --- for a simple reason: on
any $K \subseteq \cube n$ the invariant factors through $\dP(\cube n)_0^1$, which is contractible
\cite[Prop.~6.1]{PZ}.

For real hyperplane arrangements this kind of collapse is repaired by passing to the complexified
complement, whose combinatorics is recorded by Salvetti's poset. We carry that move over to
precubical sets. For any $K$ we build a category $\Chs(K)$ --- cube chains decorated by a linear
ordering of their events --- as the category of elements of a presheaf $\Lines(K)$ on $\Ch(K)$,
and show that $\Chs(\cube n) \cong \Sal(n)\op$, the Salvetti poset of the braid arrangement
$A_{n-1}$. By Salvetti's theorem $|\Chs(\cube n)| \simeq \Conf(n, \C)$, so
$\pi_1|\Chs(\cube n)| \cong P_n$ is the pure braid group --- where $|\Ch(\cube n)|$ was
contractible. Independently of that comparison, a ``no double crossing'' property of $\Chs(K)$
yields for \emph{every} $K$ a functor $\Conc : \Chs(K) \to \BraidPos$ into the positive braid
monoids in their Garside germ presentation, natural in $K$, and non-trivial as soon as $K$
contains a square. The construction, the Salvetti comparison and the braid functor are formalised
in Lean~4 \cite{Cubical}; Salvetti's theorem itself is used as a citation, not re-proved.
\end{abstract}

\tableofcontents
\bigskip

\section{Introduction}\label{sec:intro}

\subsection{The problem}

Let $K$ be a bipointed precubical set, thought of as a concurrent program: cells are states,
directions are events, and a directed path from $\init_K$ to $\fin_K$ is an execution. The
homotopy type of the execution space $\dP(K)_0^1$ is the basic invariant of the subject
\cite{FGHMR}, and Ziemiański computes it combinatorially.

\begin{theorem}[{\cite[Thm.~7.6]{Z2}, restated as \cite[Thm.~2.5]{PZ}}]\label{thm:ziemianski}
  For every bipointed precubical set $K$ the spaces $\dP(K)_0^1$ and $|\Ch(K)|$ are naturally
  weakly homotopy equivalent, as functors $\BPSet \to \mathrm{ho}\mathbf{Top}$.
\end{theorem}

Paliga and Ziemiański then show that the space of directed loops on the final precubical set $Z$
is homotopy equivalent to the total unordered configuration space of the plane
\cite[Thm.~1.1]{PZ}. Since every $K$ has a unique map $K \to Z$, this transports to invariants of
$K$: for each $\alpha \in \dP(K, n)_0^1$ a representation
\[
  IB_n(K,\alpha) : \pi_1(\dP(K,n), \alpha) \longrightarrow \pi_1(\UConf(n,\R^2)) \cong B_n,
\]
and characteristic classes pulled back from $H^*(B_n)$. And then the difficulty:

\begin{theorem}[{\cite[Prop.~6.1]{PZ}}]\label{thm:pz-trivial}
  If $K \subseteq \cube n$ then the classifying map $R_n(K)$ is null-homotopic. In particular
  $IB_n(K,\alpha)$ and every pulled-back characteristic class are trivial, for all
  $\alpha \in \dP(K,n)_0^1$.
\end{theorem}

The proof is one line: $R_n(K)$ factors through $\dP(\cube n)_0^1$, which is contractible. The
hypothesis $K \subseteq \cube n$ is exactly the condition of being a \emph{sculpture} in the sense
of Pratt \cite{Pratt}, equivalently a \emph{Euclidean cubical complex} in the sense of
Raussen--Ziemiański \cite{RZ}; the two classes agree by \cite{Z2}. That is the class most
concurrency examples live in, so the invariant is trivial precisely where one wants it.

\subsection{The proposal}

The same phenomenon is familiar for real hyperplane arrangements: the order complex of the face
poset of a real arrangement is contractible, and the fix is to complexify. The combinatorial
shadow of the complexified complement is the Salvetti poset \cite{Salvetti, Paris,
BjornerZiegler}, whose cells are a face together with a chamber above it.

This paper carries that fix over. The observation driving it is that a cube chain of $\cube n$ is
the same thing as an ordered set partition of the $n$ directions, hence the same thing as a
covector of the braid arrangement $A_{n-1}$; under this dictionary a maximally refined chain --- a
\emph{run}, all of whose beads are edges --- is a tope. So we decorate each cube chain with a run
refining it, exactly as Salvetti decorates each face with a tope above it. \Cref{def:chstar} does
this for an arbitrary precubical set, by exhibiting runs as a presheaf and taking a category of
elements.

\subsection{Results}

\begin{theorem}[\Cref{thm:sal}]\label{thm:A}
  For every $n \ge 0$ there is an isomorphism of posets $\Chs(\cube n) \cong \Sal(n)\op$, where
  $\Sal(n)$ is the Salvetti poset of the braid arrangement $A_{n-1}$. Consequently
  \[
    |\Chs(\cube n)| \;\simeq\; \Conf(n, \C),
    \qquad
    \pi_1|\Chs(\cube n)| \;\cong\; P_n,
  \]
  the pure braid group on $n$ strands, while $|\Ch(\cube n)| \simeq \dP(\cube n)_0^1$ is
  contractible.
\end{theorem}

So the complexification does not collapse on the space that made \Cref{thm:pz-trivial} work. The
braid data can also be extracted directly, without passing through Salvetti's theorem at all:

\begin{theorem}[\Cref{thm:conc}]\label{thm:B}
  For every bipointed precubical set $K$ there is a functor $\Conc : \Chs(K) \to \BraidPos$ into
  the positive braid monoids, sending a morphism to the germ of the permutation by which it
  reorders events. It is natural in $K$, and is the pullback along $\Chs(K) \to \Chs(Z)$ of a
  single universal functor on $\Chs(Z) \cong \Wdg/\rho$.
\end{theorem}

\begin{theorem}[\Cref{prop:nontrivial}]\label{thm:C}
  If there is a map $\cube 2 \to K$, then $\Conc$ is not constant at the identity braid.
\end{theorem}

\Cref{thm:B} rests on a ``no double crossing'' property (\Cref{lem:no-double-crossing}): along a
composite in $\Chs(K)$ no two events cross twice, so the associated permutations are
length-additive and multiply in the Garside germ presentation of the positive braid monoid
\cite{Garside, Garsidebook}. Compared with $IB_n$ this invariant is basepoint-free --- it is a
functor on a category rather than a conjugacy class of homomorphism out of a fundamental group ---
and it lands over the \emph{ordered} configuration space, refining the unordered picture of
\cite{PZ}; see \Cref{rem:vs-IB}, and \Cref{rem:no-section} for why no basepoint is needed and,
indeed, why none can be chosen coherently.

\subsection{Formalisation}

The constructions and combinatorial statements above --- the presheaf $\rho$, the category
$\Chs(K)$, the discrete fibration $\Chs(K) \to \Ch(K)$, the comparison of \Cref{thm:A} with the
Salvetti poset, and the braid functor of \Cref{thm:B} --- are formalised in Lean~4 on top of
mathlib \cite{Cubical}. What is \emph{not} formalised, and is used here as a citation, is the
topological input: Salvetti's theorem, \Cref{thm:ziemianski} and \Cref{thm:pz-trivial}. The
grading of \Cref{sec:next}\eqref{item:grading} is not formalised either; it is a conjecture.

\subsection{Outline}

\Cref{sec:background} recalls the cube chain material from \cite{PZ, Z2}. \Cref{sec:wedge} sets up
the Kan-extension machinery that lets us define things on cubes and inherit them on all of
$\BPSet$. \Cref{sec:arrangement} recalls the braid arrangement and the Salvetti poset.
\Cref{sec:complexification} constructs $\Chs$ and proves \Cref{thm:A}. \Cref{sec:braiding}
constructs $\Conc$ and proves \Cref{thm:B,thm:C}. \Cref{sec:next} lists what is open.

\subsection{Notation}

To fix notation, let $\Cube$ be the box category: its objects are the cubes $\cube{n}$ for
$n \ge 0$, and its morphisms are the face embeddings (composites of the coface maps
$\delta_i^\epsilon$). A \emph{precubical set} is a presheaf on $\Cube$; write $\pSet$ for the
category of precubical sets. Concretely, a precubical set $K$ has, for each $n \ge 0$, a set $K_n$
of $n$-cells, together with face maps
\[
  d_i^\epsilon : K_n \to K_{n-1} \qquad (1 \le i \le n,\ \epsilon \in \{0,1\}),
\]
where $i$ names a direction and $\epsilon$ a side, subject to the precubical identities
$d_i^\epsilon d_j^\eta = d_{j-1}^\eta d_i^\epsilon$ for $i < j$. We write $\cube{n}$ both for the
standard $n$-cube as an object of $\Cube$ and for the presheaf it represents, and we set
$[n] := \{1, \dots, n\}$, so that $[0] = \emptyset$.

We usually work with precubical sets carrying designated start and end vertices. A \emph{bipointed}
precubical set is a precubical set $K$ with two chosen vertices $\init_K, \fin_K \in K_0$,
and morphisms are required to preserve them; write $\BPSet$ for the resulting category. Each cube
$\cube{n}$ is bipointed by its initial vertex $(0,\dots,0)$ and final vertex $(1,\dots,1)$. We
write $Z$ for the final object of $\BPSet$.

\begin{convention}
Named ambient categories are set in bold ($\Cat$, $\Set$, $\Cube$, $\Braid$); constructions applied
to an object are set upright ($\Ch(K)$, $\Run(K)$, $\Sal(n)$). Functors are written with $\to$; we
reserve $\Rightarrow$ for natural transformations. The standard $n$-cube is $\cube{n}$;
the symbol $\square$ names the category of precubical sets.
\end{convention}

\begin{remark}[Variance]\label{rem:variance}
Two order conventions have to be pinned down, because they interact.

Morphisms of $\Ch(K)$ run from a \emph{finer} chain to a \emph{coarser} one (\Cref{def:ch}), while
morphisms of $\Sal(n)$ run from a smaller face to a larger one (\Cref{def:sal}). Under the
dictionary of \Cref{lem:ch-face} a finer chain is a \emph{larger} covector, so these two
conventions are opposite to each other. We define $\Chs(K)$ as the category of elements of a
presheaf on $\Ch(K)$, which makes $\pi : \Chs(K) \to \Ch(K)$ a discrete \emph{fibration} over
$\Ch(K)$; the price is the $\op$ in \Cref{thm:A}. The Lean development \cite{Cubical} makes the
opposite choice, indexing over $\Ch(K)\op$, so that its projection is a discrete
\emph{opfibration} and its comparison with $\Sal(n)$ carries no $\op$. Nothing homotopical turns on
this: $|C| \cong |C\op|$ for any small category $C$. We flag it only so that the two sources can be
read side by side.
\end{remark}

\section{Background: cube chains}\label{sec:background}

We recall three constructions on $\BPSet$ from \cite{PZ, Z2}.

The first is the \textbf{wedge} $\vee$. Given $P, Q \in \BPSet$, the wedge $P \vee Q$ is the
pushout, formed in $\pSet$, gluing $P$'s final vertex to $Q$'s initial vertex:
\[
\begin{tikzcd}
  \cube{0} \arrow[r, "\init_Q"] \arrow[d, "\fin_P"'] & Q \arrow[d] \\
  P \arrow[r] & P \vee Q \arrow[ul, phantom, very near start, "\ulcorner"]
\end{tikzcd}
\]
Note that the two legs are maps of $\pSet$, not of $\BPSet$: a bipointed map $\cube 0 \to Q$ would
have to send the point to both $\init_Q$ and $\fin_Q$. The wedge $P \vee Q$ is bipointed by
$\init_P$ and $\fin_Q$, and $(\BPSet, \vee, \cube{0})$ is monoidal with unit the point $\cube{0}$.
The \textbf{serial wedge} of a list of \emph{positive} naturals is
\[
  \bigvee[a_1, \dots, a_r] := \cube{a_1} \vee \cdots \vee \cube{a_r},
\]
with $\bigvee[\,] = \cube{0}$. Write $\Wdg \subseteq \BPSet$ for the full subcategory of serial
wedges.

\begin{definition}[{\cite[\S2]{PZ}}]\label{def:ch}
For $K \in \BPSet$, let $\Ch(K)$ be the full subcategory of the slice $\BPSet / K$ spanned by the
serial wedges: an object is a bipointed map $\bigvee a \to K$, and a morphism
$(\bigvee a \to K) \to (\bigvee b \to K)$ is a bipointed map $\bigvee a \to \bigvee b$ commuting
over $K$. Such a morphism runs from a finer chain to a coarser one.
\end{definition}

This is verbatim the cube chain category of \cite[\S2]{PZ} --- they take $\square$-maps
$\bigvee^{\vec n} \to K$ as objects and commuting triangles over $K$ as morphisms --- so
\Cref{thm:ziemianski} applies to it as stated.

The third construction is \textbf{refinement}. A \emph{chain} in $K$ is a finite sequence
$[c_1, \dots, c_l]$ of cells of $K$, each of dimension at least $1$, such that
$\fin(c_i) = \init(c_{i+1})$ for all $i$, with $\init(c_1) = \init_K$ and $\fin(c_l) = \fin_K$;
here $\init(c)$ and $\fin(c)$ denote the images under $c : \cube{\dim c} \to K$ of the initial and
final vertices of the cube. A \emph{refinement} replaces a cell $c_i$ by a chain of the cube
$\cube{\dim c_i}$, post-composed with $c_i$. Let $\Refine(K)$ be the category of chains, with
morphisms generated by refinements and oriented as in $\Ch(K)$ (finer to coarser).

Finally, two properties of $K$ from \cite{PZ} recur. We say $K$ \emph{admits altitude} if there is
a function $\alt : K_0 \to \Z$ with $\alt(\fin(e)) = \alt(\init(e)) + 1$ for every edge
$e \in K_1$. Following \cite[\S2]{PZ} verbatim, $K$ is \emph{non-self-linked} (NSL) if every
$\square$-map $\cube{n} \to K$ is injective, for all $n \ge 0$. By Yoneda a $\square$-map
$\cube n \to K$ is an $n$-cell of $K$, so this says the iterated faces of every cell are pairwise
distinct.

We use the following, all from \cite{PZ}.

\begin{lemma}[{\cite[\S2]{PZ}}]\label{lem:wedge-monoidal}
  The $\vee$ operator is a monoidal product on $\BPSet$, with unit $\cube 0$.
\end{lemma}

\begin{lemma}\label{lem:ch-lax}
  $\Ch : \BPSet \to \Cat$ is a lax monoidal functor from $(\BPSet, \vee)$ to $(\Cat, \times)$: the
  comparison $\Ch(P) \times \Ch(Q) \to \Ch(P \vee Q)$ concatenates chains.
\end{lemma}

\begin{lemma}[{\cite[Lem.~2.11]{PZ}}]\label{lem:ch-refine}
  Let $K$ admit altitude and be NSL. Then $\Ch(K)$ is a poset, every chain is determined by its set
  of cubes, and every bipointed map $\bigvee a \to K$ is injective. Consequently $\Ch(K)$ and
  $\Refine(K)$ are equivalent, and the comparison of \Cref{lem:ch-lax} is an equivalence on wedge
  decompositions of $K$.
\end{lemma}

\begin{lemma}\label{lem:wedges-nice}
  Serial wedges admit altitude and are NSL. Explicitly, $\alt$ counts edges from $\init$, so the
  final vertex of $\bigvee a$ has altitude $\sum_i a_i$, and the vertices of the $j$-th summand are
  exactly those of altitude in $[s^a_j,\ s^a_{j+1}]$, where $s^a_j := \sum_{j' < j} a_{j'}$.
\end{lemma}

\begin{lemma}[{\cite[Prop.~2.7]{PZ}}]\label{lem:alt-preserved}
  Let $P, Q$ be connected bipointed precubical sets admitting altitude functions. Then every
  bipointed map $f : P \to Q$ preserves altitude: $\alt_Q \circ f = \alt_P$.
\end{lemma}

\begin{proof}[Proof sketch]
  A precubical map takes edges to edges, so $\alt_Q$ increases by exactly one along the image of
  each edge; connectivity propagates this, and $f(\init_P) = \init_Q$ pins the constant.
\end{proof}

This is all the foundational material we need.

\section{Wedge machinery}\label{sec:wedge}

In an effort to keep our proofs better factored, we establish some results here that let us
extend constructions on $\Cube$ to constructions on $\BPSet$.

We write $E$ for both extensions below. The variance of $P$ determines which is meant: a covariant
$P$ is extended by left Kan extension, a contravariant one by right Kan extension. In both cases
the extension is taken along the Yoneda embedding $\Cube \hookrightarrow \pSet$, producing a
functor on $\pSet$, which we then restrict along the forgetful functor $\BPSet \to \pSet$ without
further comment. Since Yoneda is fully faithful, $E(P)|_{\Cube} \cong P$.

\begin{lemma}\label{lem:kan-left}
  Any functor $P : \Cube \to \Cat$ extends, by left Kan extension along $\Cube \hookrightarrow \pSet$,
  to a functor $E(P) : \BPSet \to \Cat$. If $P(\cube{0}) = \emptyset$, then $E(P)$ is strong monoidal
  \[
    E(P) : (\BPSet, \vee, \cube{0}) \to (\Cat, \sqcup, \emptyset),
  \]
  and for every serial wedge
  \[
    E(P)(\bigvee a) \;\simeq\; \coprod_i P(\cube{a_i}).
  \]
\end{lemma}
\begin{proof}
  As a pointwise left Kan extension (a coend over $\Cube$), $E(P)$ is cocontinuous on $\pSet$. The
  wedge is a pushout in $\pSet$, so $E(P)$ sends the span $A \leftarrow \cube{0} \to B$ to
  $E(P)(A \vee B) = E(P)(A) \sqcup_{P(\cube{0})} E(P)(B)$. When $P(\cube{0}) = \emptyset$ this
  pushout is the coproduct, so the comparison map is an isomorphism; the formula for $\bigvee a$
  follows by induction on the number of summands.
\end{proof}

\begin{lemma}\label{lem:kan-right}
  Any functor $P : \Cube\op \to \Cat$ extends, by right Kan extension along $\Cube \hookrightarrow \pSet$,
  to a $\Cat$-valued presheaf $E(P) : \BPSet\op \to \Cat$. If $P(\cube{0})$ is terminal, then $E(P)$
  is strong monoidal
  \[
    E(P) : (\BPSet\op, \vee, \cube{0}) \to (\Cat, \times, \mathbf 1),
  \]
  and for every serial wedge
  \[
    E(P)(\bigvee a) \;\simeq\; \prod_i P(\cube{a_i}).
  \]
  When $P$ is $\Set$-valued --- that is, when $P$ is itself a precubical set --- the extension is
  simply
  \begin{equation}\label{eq:E-is-hom}
    E(P)(X) \;=\; \Hom_{\pSet}(X, P).
  \end{equation}
\end{lemma}
\begin{proof}
  Dually, $E(P)$ is a pointwise right Kan extension, hence continuous, and being contravariant it sends the
  wedge pushout to the pullback $E(P)(A \vee B) = E(P)(A) \times_{P(\cube{0})} E(P)(B)$. When $P(\cube{0})$
  is terminal this pullback is the product, so the comparison is an isomorphism; the formula for $\bigvee a$
  follows by induction. For \eqref{eq:E-is-hom}: $\Hom_{\pSet}(-,P)$ is continuous and restricts
  along Yoneda to $P$, so it is the right Kan extension.
\end{proof}

We now instantiate this machinery. The first case is the coordinate functor.
\begin{definition}\label{def:coord}
  Let $\Coord : \Cube \to \Set$ send the cube $\cube{n}$ to its set of directions $[n]$, and send a
  face embedding $f : \cube{n} \to \cube{m}$ to the monotone injection $[n] \hookrightarrow [m]$
  that names the $n$ directions $f$ leaves free. Note $\Coord(\cube{0}) = [0] = \emptyset$, so
  \Cref{lem:kan-left} applies to $\Coord$: the extension $E(\Coord)$ is strong monoidal, with
  \[
    E(\Coord)(\bigvee a) \;\simeq\; \coprod_i [a_i].
  \]
  For a morphism $f : \bigvee a \to \bigvee b$ we write $f_* := E(\Coord)(f)$ for the induced
  function on coordinate sets,
  \[
    f_* : \coprod_i [a_i] \longrightarrow \coprod_j [b_j],
  \]
  which records where each coordinate of $\bigvee a$ is sent. We order $\coprod_i [a_i]$
  lexicographically: first by summand index, then within a summand.
\end{definition}

\begin{definition}\label{def:blk}
  Write $\blk(c)$ for the summand index of a coordinate $c \in \coprod_i [a_i]$.
  By the pushout defining $\bigvee a$, a morphism $f : \bigvee a \to \bigvee b$
  is determined by its restrictions to the summands, and each such restriction is a cell of
  $\bigvee b$ of positive dimension, hence lies in a single summand: a face embedding
  \[
    f_i : \cube{a_i} \longrightarrow \cube{b_{\blk(f)(i)}}.
  \]
  This defines $\blk(f)$, and describes $f_*$ on coordinates:
  \begin{equation}\label{eq:blk-coord}
    f_*(i, d) \;=\; \bigl(\blk(f)(i),\ (f_i)_*(d)\bigr), \qquad d \in [a_i].
  \end{equation}
  In particular $\blk \circ f_* = \blk(f) \circ \blk$.
\end{definition}

\begin{lemma}\label{lem:blk}
  For $f : \bigvee a \to \bigvee b$, the map $\blk(f)$ is monotone and $f_*$ restricts to an
  order-preserving injection $[a_i] \hookrightarrow [b_{\blk(f)(i)}]$. Consequently $\blk \circ f_*$
  is monotone.
\end{lemma}
\begin{proof}
  The restriction of $f_*$ to $[a_i]$ is $(f_i)_*$, which is an order-preserving injection because
  $f_i$ is a face embedding (\Cref{def:coord}).

  For monotonicity of $\blk(f)$, use altitude. By \Cref{lem:wedges-nice} the vertices of the $j$-th
  summand of $\bigvee b$ are exactly those of altitude in $[s^b_j,\ s^b_{j+1}]$, and these
  intervals increase with $j$; likewise the $i$-th summand of $\bigvee a$ occupies altitudes
  $[s^a_i,\ s^a_{i+1}]$. By \Cref{lem:alt-preserved} the bipointed map $f$ preserves altitude, so
  the image of summand $i$ sits at altitudes at most those of the image of summand $i+1$. Hence
  $\blk(f)(i) \le \blk(f)(i+1)$.

  The last claim follows since $\blk \circ f_* = \blk(f) \circ \blk$ and $\blk$ is monotone, the
  order on $\coprod_i[a_i]$ being lexicographic in it.
\end{proof}

\begin{lemma}\label{lem:coord-injective}
  For any morphism $f : \bigvee a \to \cube{m}$ in $\pSet$, the function $f_*$ is injective.
\end{lemma}
\begin{proof}
  We recurse on $a$.
  If $a$ is empty, then $f_* : [0] \to [m]$, which is always injective.
  For the induction case, we have $f : (\bigvee a) \vee \cube{n} \to \cube{m}$.
  Now, by the universal property of the pushout, this splits into $h : \cube{n} \to \cube{m}$
  and $g : \bigvee a \to \cube{m}$. Now, $g_*$ is injective by hypothesis.
  On the other hand, $h$ is a map of representables, and $E$ is the left Kan extension.
  So $h_* = \Coord(h)$ is the injective monotone map $[n] \to [m]$ corresponding to $h$.

  Finally, we show the two images are disjoint. Let $c \in \{0,1\}^m$ be the image of the glued vertex
  $\fin(\bigvee a) = \init(\cube{n})$.
  Maps into a cube are monotone, raising one coordinate from $0$ to $1$ per edge.
  Since $h$ starts at $c$, every point in the image of $h$ is reachable from $c$, so $h$ can only
  raise coordinates, and every direction free in $h$ is $0$ at $c$.
  On the other hand $g$ ends at $c$, so every direction $g$ raises has already reached $1$ at $c$.
  Hence
  \[
    \im h_* \subseteq \{k : c_k = 0\}
    \quad\text{and}\quad
    \im g_* \subseteq \{k : c_k = 1\}
  \]
  are disjoint. So $f_* = g_* \sqcup h_*$ is the disjoint union of injective functions with
  disjoint images, using that $E(\Coord)$ carries the defining pushout to the coproduct
  (\Cref{lem:kan-left}). Therefore $f_*$ is injective.
\end{proof}


\begin{lemma}\label{lem:coord-bijective}
  For any morphism $f : \bigvee a \to \bigvee b$ in $\BPSet$, the function $f_*$ is bijective.
\end{lemma}
\begin{proof}
  First, injectivity, by induction on the length of $b$, generalising over $a$.
  If $b$ is empty, then $\bigvee b = \cube{0}$, which has no cells above dimension $0$; so
  $\bigvee a$ has none either, forcing $a$ to be empty and $f_* : [0] \to [0]$.

  Now consider $f : \bigvee a \to \cube{m} \vee \bigvee b$. The inductive hypothesis is that for
  any $x$ the coordinate map of a morphism $\bigvee x \to \bigvee b$ is injective. By
  \Cref{lem:wedges-nice} the target admits altitude and is NSL, so by \Cref{lem:ch-refine} the
  comparison $\Ch(\cube{m}) \times \Ch(\bigvee b) \to \Ch(\cube{m} \vee \bigvee b)$ is an
  equivalence. Viewing $f$ as a morphism of $\Ch(\cube{m} \vee \bigvee b)$ from
  $(\bigvee a \to \cube m \vee \bigvee b)$ to the identity object, it therefore splits as a pair
  $(f_1, f_2)$: the list $a$ splits as a concatenation $a' \ast a''$ with
  $f_1 : \bigvee a' \to \cube{m}$ and $f_2 : \bigvee a'' \to \bigvee b$. Now $(f_2)_*$ is injective
  by the induction hypothesis and $(f_1)_*$ is injective by \Cref{lem:coord-injective}.
  Since $f \cong f_1 \vee f_2$ and $E(\Coord)$ is strong monoidal,
  \[
    f_* = (f_1 \vee f_2)_* = (f_1)_* \sqcup (f_2)_*,
  \]
  the disjoint union of injections with disjoint codomains. Therefore $f_*$ is injective.

  For surjectivity we count. By \Cref{lem:wedges-nice} the final vertex of $\bigvee a$ has altitude
  $\sum_i a_i$, and likewise for $b$; by \Cref{lem:alt-preserved} the bipointed map $f$ preserves
  altitude and sends $\fin$ to $\fin$, so $\sum_i a_i = \sum_j b_j$. Thus $f_*$ is a
  cardinality-preserving injection of finite sets, hence bijective.
\end{proof}

\section{The braid arrangement and the Salvetti poset}\label{sec:arrangement}

The construction to come is a transcription of the following, so we recall it in the form we will
imitate. We define the braid arrangement and its related constructions following Paris
\cite{Paris}, specialised to the braid arrangement for convenience.

The braid arrangement $A_{n-1}$ is the family of hyperplanes
\[
  H_{ij} = \{ x \in \R^n : x_i = x_j \}, \qquad 1 \leq i < j \leq n,
\]
one for each pair of coordinates. Orient $H_{ij}$ by the sign of $x_i - x_j$. Each point $x \in \R^n$ then
assigns to every hyperplane a sign in $\{-,0,+\}$, and the sign vectors so obtained are the
\textbf{covectors}, or faces, of the arrangement:
\[
  \Face(n) := \bigl\{\, F : \{(i,j) : 1 \leq i < j \leq n\} \to \{-,0,+\} \ \bigm|\ F \text{ arises from some } x \in \R^n \,\bigr\}.
\]
A covector records the weak order of the coordinates of $x$, so it is the same data as an \emph{ordered
set partition} of $[n]$: the blocks are the coordinates of equal value, listed in increasing
order of that value. We order covectors by
\[
  F \leq G \quad:\!\iff\quad \forall h,\ F(h) = 0 \ \text{ or }\ F(h) = G(h),
\]
so that $G$ resolves some of the ties left open by $F$. The maximal covectors --- those never taking the
value $0$ --- are the \textbf{topes}. A tope has all coordinates distinct, hence is a total order on
$[n]$; there are exactly $n!$ of them, one per permutation.

Sadly, the order complex of $\Face(n)$ is contractible --- this is the arrangement-theoretic form
of the contractibility that powers \Cref{thm:pz-trivial}. Salvetti's construction reruns the story
over $\C$ and finds much richer homotopy.

\begin{definition}\label{def:sal}
The \textbf{Salvetti poset} is
\[
  \Sal(n) := \bigl\{\, (F, C) \ \bigm|\ F \in \Face(n),\ C \text{ a tope},\ F \leq C \,\bigr\},
\]
whose order uses the composition of covectors,
\[
  (F \circ G)(h) := \begin{cases} F(h) & \text{if } F(h) \neq 0, \\ G(h) & \text{if } F(h) = 0, \end{cases}
\]
namely
\[
  (F, C) \preceq (F', C') \quad:\!\iff\quad F \leq F' \ \text{ and }\ C' = F' \circ C.
\]
\end{definition}

This is the order of \cite{Paris}; see \cite{DDP} for the form used in the formalisation.

\begin{theorem}[Salvetti \cite{Salvetti}; see also \cite{BjornerZiegler, Paris}]\label{thm:salvetti}
  $\Sal(n)$ is the face poset of a regular CW complex homotopy equivalent to the complement of the
  complexified braid arrangement in $\C^n$, namely the ordered configuration space
  $\Conf(n,\C) = \{z \in \C^n : z_i \neq z_j \text{ for } i \neq j\}$. Since the order complex of
  the face poset of a regular CW complex is its barycentric subdivision, $|\Sal(n)| \simeq
  \Conf(n,\C)$ and $\pi_1|\Sal(n)| \cong P_n$, the pure braid group.
\end{theorem}

The packaging we will imitate is as a Grothendieck construction. Let
\[
  \Tope : \Face(n) \to \Set, \qquad \Tope(F) := \{ C \text{ a tope} : F \leq C \},
\]
be the covariant functor sending $F \leq F'$ to the map $C \mapsto F' \circ C$ (composing on the left by
$F'$ carries a tope above $F$ to one above $F'$). Then
\[
  \Sal(n) \;\cong\; \textstyle\int \Tope,
\]
the Grothendieck construction: its objects are the pairs $(F, C)$ with $C \in \Tope(F)$, and its morphisms
are exactly the relations $\preceq$ above. Intuitively, each face is annotated with a compatible
tope, a local orientation; crossing a face can flip the orientation, and these flips are the braid
letters. The construction of \cite{PZ} retains the same orientation data, breaking every tie at a
face in all possible ways.

\section{The complexification \texorpdfstring{$\Chs$}{Ch*}}\label{sec:complexification}

The bridge between the two sides is a way of reading a cube chain as execution data. Associate to
each coordinate of the $n$-cube a real number: the time at which that event occurs. Then the
coordinates of $\R^n$ behave like execution information --- when did each event occur --- and the
regions cut out by the equalities $x_i = x_j$ correspond to wedge maps: sets of equal coordinates
give coarser wedge maps. That is exactly the braid arrangement of \Cref{sec:arrangement}. To
rebuild it generically on the precubical side, we start with the corresponding definition of tope.

\subsection{Runs}

\begin{definition}
  A serial wedge is \textbf{one dimensional} if all its entries are $1$.
\end{definition}

\begin{definition}
  Let $\Run(K)$ be the full subcategory of $\Ch(K)$ whose objects have
  one dimensional serial wedges as domain. We write $\Run(K)$ also for its set of objects when no
  confusion arises.
\end{definition}

\begin{lemma}\label{lem:run-monoidal}
  Being one dimensional is preserved and reflected by $\vee$. Hence on serial wedges the
  comparison of \Cref{lem:ch-refine} restricts to a bijection
  $\Run(\bigvee a) \cong \prod_i \Run(\cube{a_i})$.
\end{lemma}

Now, the morphisms we want to build are a bit unusual, because they are built from
degeneracies, which are not morphisms of $\pSet$. That is, we want a contravariant assignment
from $\Ch(K)$ to runs, mirroring the order from Salvetti:
\[
  (\bigvee a \to \bigvee b) \;\longmapsto\; (\Run(\bigvee b) \to \Run(\bigvee a)).
\]

To achieve this we use the cell-list presentation of a run. Let $\Edge(W)$ be the set of chains in
$\Refine(W)$ all of whose cells are one dimensional. Serial wedges are NSL and admit altitude
(\Cref{lem:wedges-nice}), so $\Ch(W) \simeq \Refine(W)$ there by \Cref{lem:ch-refine}, and
restricting that equivalence to the one dimensional part gives a bijection
\[
  \Run(W) \;\cong\; \Edge(W)
\]
for every serial wedge $W$: a one dimensional wedge over $W$ is the same data as its list of edges.

Now, given a face embedding $f : \cube{n} \to \cube{m}$, it defines a unique
projection of cells of $\cube{m}$ to cells of $\cube{n}$, non-increasing on dimension: it forgets the
directions of $\cube{m}$ that $f$ does not leave free. Call it $\hat{f}$. So we define
\[
  \rho(f) : \Edge(\cube{m}) \to \Edge(\cube{n}), \qquad
  [c_1,\dots,c_k] \mapsto \normalize[\hat{f}(c_1), \dots, \hat{f}(c_k)],
\]
where $\normalize$ removes all the $0$-cells, and transport it along $\Run \cong \Edge$ to get
\[
  \rho(f) : \Run(\cube{m}) \to \Run(\cube{n}).
\]
Concretely, a run of $\cube{m}$ is a bijection $r_* : [m] \to [m]$ (\Cref{lem:coord-bijective}), and
$\rho(f)$ restricts it along $f_*$: the run $\rho(f)(r)$ is the unique one satisfying
\begin{equation}\label{eq:rho-restrict}
  f_* \circ \rho(f)(r)_* \;=\; r_* \circ \iota
  \qquad\text{for some strictly increasing } \iota : [n] \to [m],
\end{equation}
that is, it visits the directions left free by $f$ in the order $r$ meets them. Uniqueness holds
because $\im(f_* \circ \rho(f)(r)_*) = \im f_*$ forces $\im\iota = r_*^{-1}(\im f_*)$, which
determines $\iota$ and hence $\rho(f)(r)_* = f_*^{-1} \circ r_* \circ \iota$.

Note that the result is indeed a one dimensional chain in $\cube{n}$ because
\begin{enumerate}
\item $\hat{f}$ never increases dimension, so all dimensions are less than or equal to $1$;
\item normalization removes all the $0$-cells, so every remaining cell is one dimensional;
\item every projection respects initial and final points, so consecutive cells still match up and
      the endpoints are still $\init$ and $\fin$.
\end{enumerate}

\begin{lemma}\label{lem:rho-functor}
  $\rho : \Cube\op \to \Set$ is a functor, and $\rho(\cube{0})$ is a singleton. In particular
  $\rho$ is a precubical set, with $\rho(\cube n) \cong \Sigma_n$ and $d_i^0 = d_i^1$.
\end{lemma}
\begin{proof}
  Cubes are NSL and admit altitude, so $\Run(\cube{n}) \cong \Edge(\cube{n})$ as sets and it is
  enough to work with cell lists. The projection attached to an identity embedding is the identity,
  so $\rho$ preserves identities. If $f : \cube{n} \to \cube{m}$ and $g : \cube{m} \to \cube{l}$
  are face embeddings then $\widehat{g \circ f} = \hat f \circ \hat g$, since both forget exactly
  the directions not left free by $g \circ f$. It remains to see that normalising once at the end
  agrees with normalising after each step: $\hat f$ carries $0$-cells to $0$-cells, so deleting
  them early or late deletes the same cells. Hence $\rho(g \circ f) = \rho(f) \circ \rho(g)$.

  For the singleton, observe that $\rho(\cube{0}) = \Run(\cube{0})$,
  but the only one dimensional serial wedge which maps into $\cube{0}$ is the empty one.
  So indeed $\rho(\cube{0})$ is a singleton. The description of $\rho(\cube n)$ is
  \Cref{lem:coord-bijective}: a run of $\cube n$ is a bijection $[n] \to [n]$. The face map
  $\rho(\delta_i^\epsilon)$ deletes direction $i$ from the permutation, visibly independently of
  $\epsilon$.
\end{proof}

That means two things. Firstly, $\rho$ is itself a precubical set, so \eqref{eq:E-is-hom} applies.
Second, our lifting machinery sends $\rho$ to $E(\rho) : \BPSet\op \to \Set$. The next lemma is the
linchpin: it says $\rho$ \emph{classifies} runs.

\begin{lemma}[$\rho$ classifies runs]\label{lem:rho-classifies}
  For every serial wedge $\bigvee a$ there is a bijection
  \[
    E(\rho)\bigl(\textstyle\bigvee a\bigr) \;=\; \Hom_{\pSet}\bigl(\textstyle\bigvee a,\ \rho\bigr)
      \;\cong\; \Run\bigl(\textstyle\bigvee a\bigr),
  \]
  natural in $\bigvee a$, under which
  \begin{enumerate}
    \item a run $x$ corresponds to the tuple $(x_1, \dots, x_r)$ of its restrictions
          $x_i \in \Run(\cube{a_i})$ to the summands; and
    \item for a morphism $f : \bigvee a \to \bigvee b$ of $\BPSet$ and $y \in \Run(\bigvee b)$, the
          $i$-th component of $E(\rho)(f)(y)$ is
          \begin{equation}\label{eq:rho-componentwise}
            \bigl(E(\rho)(f)(y)\bigr)_i \;=\; \rho(f_i)\bigl(y_{\blk(f)(i)}\bigr).
          \end{equation}
  \end{enumerate}
\end{lemma}
\begin{proof}
  By \Cref{lem:rho-functor} $\rho(\cube 0)$ is terminal, so \Cref{lem:kan-right} applies and, using
  \eqref{eq:E-is-hom} and Yoneda,
  \[
    \Hom_{\pSet}\bigl(\textstyle\bigvee a, \rho\bigr) \;\cong\; \prod_i \Hom_{\pSet}(\cube{a_i}, \rho)
      \;=\; \prod_i \rho(\cube{a_i}) \;=\; \prod_i \Run(\cube{a_i});
  \]
  concretely, a map out of the pushout is a compatible family of maps out of the summands, and the
  compatibility conditions live over $\rho(\cube 0)$, a singleton, hence are vacuous. On the other
  side $\Run(\bigvee a) \cong \prod_i \Run(\cube{a_i})$ by \Cref{lem:run-monoidal}. Composing gives
  the bijection, and by construction it is summandwise restriction, which is (1).

  For (2), $E(\rho)(f)$ is precomposition with $f$ by \eqref{eq:E-is-hom}, and the $i$-th component
  of the product decomposition is restriction along the summand inclusion
  $\cube{a_i} \hookrightarrow \bigvee a$. The composite
  $\cube{a_i} \hookrightarrow \bigvee a \xrightarrow{f} \bigvee b$ is, by \Cref{def:blk}, the face
  embedding $f_i$ followed by the inclusion of the summand $\cube{b_{\blk(f)(i)}}$. Precomposition
  with $f_i$ is $\rho(f_i)$ by definition of $\rho$, and restricting along the summand inclusion of
  $\bigvee b$ picks out $y_{\blk(f)(i)}$.
\end{proof}

From now on we silently identify $E(\rho)(\bigvee a)$ with $\Run(\bigvee a)$.

\subsection{The construction}

\begin{definition}\label{def:chstar}
Let $W(K) : \Ch(K) \to \Wdg$ be the forgetful functor $(\bigvee a \to K) \mapsto \bigvee a$. Set
\[
  \Lines(K) := E(\rho) \circ W(K)\op,
\]
a presheaf on $\Ch(K)$, and
\[
  \Chs(K) := \textstyle\int \Lines(K),
\]
the category of elements of $\Lines(K)$, with projection $\pi : \Chs(K) \to \Ch(K)$.
\end{definition}

Defining $\Chs$ abstractly has some conveniences later.
Let us unpack the definition of $\Chs$ to have a more concrete presentation of it.

Objects of $\Chs(K)$ are pairs
\[
  (\bigvee a \to K,\ x \in \Run(\bigvee a)), \text{ i.e. } \bigvee[1,\dots,1] \to \bigvee a \to K,
\]
a cube chain together with a linear ordering of its events.
The intuition here matches the Salvetti story precisely: we have a wedge map into $K$, coupled with
orientation info about it. Morphisms are a little uglier. A morphism
\[
  (\bigvee a \to K,\ x \in \Run(\bigvee a)) \longrightarrow (\bigvee b \to K,\ y \in \Run(\bigvee b))
\]
unpacks to
\begin{enumerate}
  \item A morphism $f : (\bigvee a \to K) \to (\bigvee b \to K)$,
  i.e.\ a bipointed map $f' : \bigvee a \to \bigvee b$ commuting over $K$.
  \item The condition $x = \Lines(f)(y)$:
  we take the run $y$, project it along $f'$, and require the result to be $x$. Since the fibers of
  $E(\rho)$ are discrete, this is a condition, not extra data --- $x$ is determined by $f'$ and $y$.
\end{enumerate}
In particular $\pi$ is a discrete fibration; see \Cref{rem:variance}.

The diagram below records the story. The solid arrows are morphisms of $\Chs(K)$; the dashed
$\sigma(f)$ (a permutation of positions) and $E(\rho)(f)$ (the run projection) are plain functions, not
morphisms in $\BPSet$, and they point oppositely --- $E(\rho)(f)$ carries $y$ back to $x$, and
$\sigma(f)$ is the induced reindexing.
\[
\begin{tikzcd}[row sep=huge, column sep=huge]
  \bigvee[1,\dots,1] \arrow[r, "x"] \arrow[d, perm, "\sigma(f)"']
    & \bigvee a \arrow[d, "f"] \arrow[rd] & \\
  \bigvee[1,\dots,1] \arrow[r, "y"']
    & \bigvee b \arrow[u, runact, bend left=85, "E(\rho)(f)"] \arrow[r] & K
\end{tikzcd}
\]

\subsection{The crossing permutation}

Let's look more closely at the construction of $\sigma$.
Note that a run $r \in \Run(\bigvee p)$ is itself a bipointed map
$\bigvee[1,\dots,1] \to \bigvee p$, so \Cref{def:coord} already applies to it, and by
\Cref{lem:coord-bijective} its coordinate map
\[
  r_* : [N] \xrightarrow{\ \sim\ } \textstyle\coprod_i [p_i], \qquad N := \textstyle\sum_i p_i,
\]
is exactly a linear ordering of the coordinates of $\bigvee p$.
Conjugating the coordinate bijection $f'_* : \coprod_i [a_i] \to \coprod_j [b_j]$ by the two run
orderings gives a permutation of $[N]$,
\[
  \sigma(f) := y_*^{-1} \circ f'_* \circ x_* \ \in\ \Sigma_N,
\]
equivalently $y_* \circ \sigma(f) = f'_* \circ x_*$. This is well posed: by
\Cref{lem:coord-bijective} the source and target of $f$ have the same number of coordinates.
This permutation $\sigma(f)$ is what will power our braid functor.

\begin{lemma}\label{lem:sigma-functor}
$\sigma : \Chs(K) \to \bigsqcup_N \Sigma_N$ is a functor, where the target is the groupoid with
objects the naturals and $\Hom(N,N) = \Sigma_N$.
\end{lemma}
\begin{proof}
Along a composite the middle ordering $y_*$ cancels,
\[
  \sigma(g \circ f) = z_*^{-1} \circ g'_* \circ f'_* \circ x_*
    = \bigl(z_*^{-1} \circ g'_* \circ y_*\bigr)\bigl(y_*^{-1} \circ f'_* \circ x_*\bigr)
    = \sigma(g)\,\sigma(f),
\]
and
\[
  \sigma(\id) = x_*^{-1} \circ (\id)_* \circ x_* = \id. \qedhere
\]
\end{proof}

\begin{lemma}\label{lem:within-summand}
  Let $f : (\bigvee a, x) \to (\bigvee b, y)$ in $\Chs(K)$ and suppose
  $\blk(x_*(k)) = \blk(x_*(l))$. Then $k < l$ if and only if $\sigma(f)(k) < \sigma(f)(l)$.
\end{lemma}
\begin{proof}
  By \Cref{lem:rho-classifies}, $x$ decomposes as $(x_1, \dots, x_r)$ with
  $(x_i)_* : [a_i] \to [a_i]$. Since $\blk \circ x_*$ is monotone (\Cref{lem:blk}) and $x_*$ is
  bijective (\Cref{lem:coord-bijective}), the preimage under $x_*$ of the $i$-th summand is a run of
  $a_i$ consecutive positions starting at $s^a_i := \sum_{i' < i} a_{i'}$, and
  \begin{equation}\label{eq:run-block}
    x_*(s^a_i + t) \;=\; \bigl(i,\ (x_i)_*(t)\bigr), \qquad t \in [a_i].
  \end{equation}

  Put $i := \blk(x_*(k)) = \blk(x_*(l))$ and $j := \blk(f')(i)$, and write $k = s^a_i + t$ and
  $l = s^a_i + t'$. By \eqref{eq:rho-componentwise} the morphism condition $x = E(\rho)(f')(y)$
  reads componentwise as $x_i = \rho(f_i)(y_j)$, so \eqref{eq:rho-restrict} gives
  $(f_i)_* \circ (x_i)_* = (y_j)_* \circ \iota$ with $\iota$ strictly increasing. Hence
  \[
    y_*(\sigma(f)(k))
      = f'_*\bigl(x_*(s^a_i + t)\bigr)
      = f'_*\bigl(i, (x_i)_*(t)\bigr)
      = \bigl(j, (f_i)_*((x_i)_*(t))\bigr)
      = \bigl(j, (y_j)_*(\iota(t))\bigr)
      = y_*\bigl(s^b_j + \iota(t)\bigr),
  \]
  using \eqref{eq:blk-coord} and then \eqref{eq:run-block} for $y$.
  As $y_*$ is injective, $\sigma(f)(k) = s^b_j + \iota(t)$ and $\sigma(f)(l) = s^b_j + \iota(t')$; since $\iota$
  is strictly increasing, $\sigma(f)(k) < \sigma(f)(l)$ iff $t < t'$ iff $k < l$.
\end{proof}

\subsection{Comparison with the Salvetti poset}

\begin{lemma}\label{lem:ch-face}
  There is an isomorphism of posets $\Ch(\cube n) \cong \Face(n)\op$.
\end{lemma}
\begin{proof}
  An object of $\Ch(\cube n)$ is a bipointed map $\bigvee a \to \cube n$. By
  \Cref{lem:coord-bijective} its coordinate map is a bijection $\coprod_i [a_i] \to [n]$, and by
  \Cref{lem:blk} it is order-preserving on each summand; so the object determines the ordered set
  partition of $[n]$ whose $i$-th block is the image of $[a_i]$. Conversely every ordered set
  partition arises this way, and by \Cref{lem:ch-refine} --- cubes are NSL and admit altitude ---
  the object is determined by it. Ordered set partitions of $[n]$ are exactly the covectors of
  $\Face(n)$.

  For the order: a morphism $\bigvee a \to \bigvee b$ over $\cube n$ exists exactly when the
  partition attached to $a$ refines that attached to $b$, i.e.\ when the covector of $b$ resolves
  fewer ties than that of $a$, i.e.\ when $F_b \le F_a$. Since morphisms of $\Ch$ run finer to
  coarser, this is the opposite order.
\end{proof}

\begin{theorem}\label{thm:sal}
  For every $n \ge 0$ there is an isomorphism of posets $\Chs(\cube n) \cong \Sal(n)\op$.
  Consequently $|\Chs(\cube n)| \simeq \Conf(n,\C)$ and $\pi_1|\Chs(\cube n)| \cong P_n$.
\end{theorem}
\begin{proof}
  Transport $\Lines(\cube n)$ across \Cref{lem:ch-face}. A presheaf on $\Ch(\cube n)$ is a covariant
  functor on $\Ch(\cube n)\op \cong \Face(n)$. On objects, $\Lines(\cube n)$ sends a chain to its
  set of runs; under the dictionary a run of $\bigvee a$ is a linear order on $[n]$ refining the
  partition, that is a tope $C \ge F_a$. So $\Lines(\cube n)$ and $\Tope$ agree on objects. On
  morphisms, $\Lines$ acts by projecting a run along the chain map, which by
  \eqref{eq:rho-restrict} reads off the order in which the finer chain meets the events; on
  covectors that is exactly $C \mapsto F' \circ C$, the action of $\Tope$. Hence
  $\Lines(\cube n) \cong \Tope$ as functors on $\Face(n)$.

  Now compare categories of elements. A morphism $(c, x) \to (c', y)$ of $\Chs(\cube n)$ lies over
  $c \to c'$ in $\Ch(\cube n)$, i.e.\ over $F_{c'} \le F_c$ in $\Face(n)$, and asserts
  $x = \Tope(F_{c'} \le F_c)(y)$ --- which is precisely a morphism $(F_{c'}, y) \to (F_c, x)$ of
  $\Sal(n) = \int \Tope$. So the two categories are isomorphic after reversing arrows.

  The topological statement is \Cref{thm:salvetti} together with $|C| \cong |C\op|$.
\end{proof}

The routine verifications behind \Cref{lem:ch-face,thm:sal} --- that the dictionary is a bijection
and that the two actions on topes agree --- are machine-checked in \cite{Cubical}.

\begin{remark}\label{rem:noncollapse}
  By \Cref{thm:ziemianski}, $|\Ch(\cube n)| \simeq \dP(\cube n)_0^1$ is contractible; that
  contractibility is the whole proof of \Cref{thm:pz-trivial}. \Cref{thm:sal} says the
  complexification of the same object is a $K(P_n,1)$. This is the sense in which the construction
  answers \Cref{thm:pz-trivial}.
\end{remark}

\section{The braiding}\label{sec:braiding}

You'll notice that the $K$ side is not involved in the computation of the permutation,
or the proofs of its functoriality. This comes from a rather nice fact about Grothendieck
constructions.

\begin{lemma}\label{lem:pullback}
  The square
  \[
  \begin{tikzcd}[column sep=large, row sep=large]
    \Chs(K) = {\textstyle\int} (E(\rho)\circ W(K)\op)
      \arrow[r, "Q"] \arrow[d, "\pi"']
      \arrow[dr, phantom, very near start, "\lrcorner"]
      & {\textstyle\int} E(\rho)|_{\Wdg} = \Wdg / \rho \arrow[d, "\pi"] \\
    \Ch(K) \arrow[r, "W(K)"'] & \Wdg
  \end{tikzcd}
  \]
  is a strict pullback in $\Cat$. Here $\Wdg/\rho$ is the comma category of serial wedges over
  $\rho$, formed along $\Wdg \subseteq \BPSet \to \pSet$; note $\rho$ itself is not a serial wedge.
\end{lemma}
\begin{proof}
  By definition $\Lines(K)$ is the restriction of $E(\rho)|_{\Wdg}$ along $W(K)$, and the category
  of elements of a restricted presheaf is the strict pullback of the category of elements along the
  restricting functor: both sides have as objects the pairs
  $(c \in \Ch(K),\ x \in E(\rho)(W(K)c))$, with the same morphisms. The identification
  $\int E(\rho)|_{\Wdg} = \Wdg/\rho$ is \eqref{eq:E-is-hom}: the category of elements of
  $\Hom_{\pSet}(-,\rho)$ is the comma category.
\end{proof}

The map $Q$ is just a projection to the first arrow of $\bigvee[1,\dots,1] \to \bigvee a \to K$.
And $\rho$ is the classifier for runs (\Cref{lem:rho-classifies}). Since $E$ is, in the
contravariant case, the restricted Yoneda embedding \eqref{eq:E-is-hom}, the right-hand corner is a
comma category of wedges over $\rho$. Note that $K$ does not appear on the right hand side at all.

\begin{corollary}\label{cor:functorial}
  $\Chs$ is a functor $\BPSet \to \Cat$, and the square of \Cref{lem:pullback} identifies $Q$ with
  $\Chs(!_K) : \Chs(K) \to \Chs(Z)$, where $!_K : K \to Z$ is the unique map to the final object.
\end{corollary}
\begin{proof}
  For $u : K \to L$ we have $W(L) \circ \Ch(u) = W(K)$, hence
  $\Lines(L) \circ \Ch(u)\op = \Lines(K)$, which induces $\Chs(u)$ on categories of elements.
  Taking $L = Z$ and applying \Cref{lem:chs-Z} gives the identification.
\end{proof}

\begin{lemma}\label{lem:Q-faithful}
  The functor $Q$ is always faithful.
\end{lemma}
\begin{proof}
  Let $\bigvee[1,\dots,1] \to \bigvee a \to K$ and $\bigvee[1,\dots,1] \to \bigvee b \to K$
  be two objects of $\Chs(K)$. Two morphisms $f, g$ between these objects are equal exactly when they
  have the same underlying map $\bigvee a \to \bigvee b$,
  which happens if and only if $Q(f) = Q(g)$. So $Q$ is indeed faithful.
\end{proof}
But, we can do a nicer trick.
\begin{lemma}\label{lem:chs-Z}
  Let $Z$ be the final precubical set. Then $\Chs(Z)$ is isomorphic to $\Wdg / \rho$.
\end{lemma}
\begin{proof}
  We know $Q$ is faithful already.
  To show $Q$ is full, again we take two objects $\bigvee[1,\dots,1] \to \bigvee a \to Z$ and
  $\bigvee[1,\dots,1] \to \bigvee b \to Z$.
  Let $f$ be a morphism in $\Wdg / \rho$. That is, $f : \bigvee a \to \bigvee b$ forming a commuting
  triangle:
  \[
  \begin{tikzcd}
    \bigvee a \arrow[r] \arrow[d, "f"'] & \rho \\
    \bigvee b \arrow[ur] &
  \end{tikzcd}
  \]
  But, since $Z$ is the final precubical set, every map commutes in the triangle
  \[
  \begin{tikzcd}
    \bigvee a \arrow[r] \arrow[d, "f"'] & Z \\
    \bigvee b \arrow[ur] &
  \end{tikzcd}
  \]
 So $Q$ is full.

 On objects, every wedge map $\bigvee a \to \rho$ is the same data as a
 $\bigvee[1,\dots,1] \to \bigvee a$ by \Cref{lem:rho-classifies}, and every $\bigvee a$
 has a unique map $\bigvee a \to Z$. So $Q$ is a bijection on objects, and hence an isomorphism.
\end{proof}

\subsection{Braid words from permutations}

But we need a device for building braid words from permutations, which brings us to the Garside
germ presentation of braids \cite{Garside, Garsidebook}. For $\sigma \in \Sigma_N$ write
\[
  \Inv(\sigma) := \{\, (i,j) \in [N] \times [N] \mid i < j \text{ and } \sigma(j) < \sigma(i) \,\}
\]
for its set of \emph{inversions}, and $\ell(\sigma) := |\Inv(\sigma)|$ for its \emph{length}. We
use the same notation for a bijection between two finite linearly ordered sets, the inversions
being taken with respect to those orders; the only case needed is a coordinate map
$f_* : \coprod_i [a_i] \to \coprod_j [b_j]$ with the lexicographic orders of \Cref{def:coord}.

\begin{definition}\label{def:braid}
  Let $\BraidPos$ be the category with objects the naturals and $\Hom(N,N) = B_N^+$ the positive
  braid monoid on $N$ strands in its \emph{germ presentation}: one generator $[\sigma]$ for each
  $\sigma \in \Sigma_N$, subject to
  \[
    [\sigma][\sigma'] = [\sigma\sigma'] \qquad\text{whenever}\qquad \ell(\sigma\sigma') = \ell(\sigma) + \ell(\sigma').
  \]
  Write $\Braid$ for the corresponding category of braid \emph{groups}, $\Hom(N,N) = B_N$; the
  canonical functor $\BraidPos \to \Braid$ is faithful by Garside's theorem \cite{Garside}.
\end{definition}

The key fact is a ``no double crossings'' feature of $\Chs$: along a composite, no two events cross
each other twice. Let's first state this in pure combinatorics.

\begin{lemma}\label{lem:no-double-crossing}
  Let $f$ and $g$ be permutations of $[N]$,
  and suppose that $\Inv(f) \subseteq \Inv(g \circ f)$.
  Then $\ell(g \circ f) = \ell(g) + \ell(f)$.
\end{lemma}
\begin{proof}
Every element of $\Inv(g\circ f)$ either lies in $\Inv(f)$ or not, so
\[
  \Inv(g\circ f) \;=\; \big(\Inv(g\circ f)\cap \Inv(f)\big)\ \sqcup\ \big(\Inv(g\circ f)\setminus \Inv(f)\big).
\]
By hypothesis $\Inv(f)\subseteq \Inv(g\circ f)$, so the first set equals $\Inv(f)$ and contributes $\ell(f)$.

For the second, a pair $(k,l)\notin \Inv(f)$ has $f(k)<f(l)$, and then $(k,l)\in \Inv(g\circ f)$
iff $g(f(k))>g(f(l))$, i.e. iff $(f(k),f(l))\in \Inv(g)$. Since $f$ is a bijection,
$(k,l)\mapsto(f(k),f(l))$ restricts to a bijection
\[
  \Inv(g\circ f)\setminus \Inv(f)\ \xrightarrow{\ \sim\ }\ \{(p,q)\in \Inv(g)\;:\;f^{-1}(p)<f^{-1}(q)\}.
\]
This target is all of $\Inv(g)$: otherwise take $(p,q)\in \Inv(g)$ and set
$k=f^{-1}(q)<l=f^{-1}(p)$; then $f(k)=q>p=f(l)$ gives $(k,l)\in \Inv(f)$, while
$g(f(k))=g(q)<g(p)=g(f(l))$ gives $(k,l)\notin \Inv(g\circ f)$, contradicting $\Inv(f)\subseteq \Inv(g\circ f)$.
Hence the second set contributes $\ell(g)$, and $\ell(g\circ f)=\ell(f)+\ell(g)$.
\end{proof}

\begin{theorem}\label{thm:conc}
The map $f \mapsto [\sigma(f)]$ yields a functor $\Conc : \Chs(K) \to \BraidPos$, natural in $K$.
\end{theorem}
\begin{proof}
An object goes to the braid object on its number of coordinates, well defined by
\Cref{lem:coord-bijective}, and $\sigma(\id) = \id$. For
composition we need $[\sigma(f)][\sigma(g)] = [\sigma(f) \circ \sigma(g)]$, which by the germ
presentation holds once $\ell(\sigma(f) \circ \sigma(g)) = \ell(\sigma(f)) + \ell(\sigma(g))$, and by
\Cref{lem:no-double-crossing} it is enough to prove
$\Inv(\sigma(g)) \subseteq \Inv(\sigma(f) \circ \sigma(g))$
for composable $g : (\bigvee c, r) \to (\bigvee b, q)$ and $f : (\bigvee b, q) \to (\bigvee a, p)$.

Let $(k,l) \in \Inv(\sigma(g))$ and put $u := \sigma(g)(k)$, $v := \sigma(g)(l)$, so $k < l$ and
$v < u$. Since $q_* \circ \sigma(g) = g'_* \circ r_*$, we have $q_*(u) = g'_*(r_*(k))$ and
$q_*(v) = g'_*(r_*(l))$, so
\[
  \blk(q_*(u)) = \blk(g')\bigl(\blk(r_*(k))\bigr) \;\le\; \blk(g')\bigl(\blk(r_*(l))\bigr) = \blk(q_*(v)),
\]
using $k < l$ and \Cref{lem:blk} twice. On the other hand $v < u$ gives
$\blk(q_*(v)) \le \blk(q_*(u))$. So $\blk(q_*(u)) = \blk(q_*(v))$: the two positions lie in a common
summand of $\bigvee b$. Now \Cref{lem:within-summand} applied to $f$ turns $v < u$ into
$\sigma(f)(v) < \sigma(f)(u)$, which is exactly $(k,l) \in \Inv(\sigma(f) \circ \sigma(g))$.

For naturality in $K$, note that $\sigma$ reads only the underlying wedge data, so
$\Conc_K = \Conc_Z \circ \Chs(!_K)$ by \Cref{cor:functorial}.
\end{proof}

This looks a bit silly, because all we've done is proven our braid words have no cancellation.
But that only means we found a nice basis, and in fact it means we're working in the positive braid
monoid, where by Garside's theorem \cite{Garside} the map $B_N^+ \hookrightarrow B_N$ loses no
information. We recover the full braiding structure by lifting to the free groupoid.

\begin{definition}\label{def:conc-cat}
  Let $\Free(\Chs(K))$ be the free groupoid on $\Chs(K)$ --- its localisation at all morphisms ---
  and let
  \[
    \Conc(K) : \Free(\Chs(K)) \longrightarrow \Braid
  \]
  be the unique extension of \Cref{thm:conc} along $\Chs(K) \to \Free(\Chs(K))$, which exists
  because $\Braid$ is a groupoid. We call $\Conc(K)$ --- or, once a component is fixed, the induced
  homomorphism on vertex groups --- the \emph{concurrency braid} of $K$.
\end{definition}

\begin{proposition}\label{prop:nontrivial}
For any precubical set $K$, if there is an $f : \cube{2} \to K$, then
$\Conc(K)$ is nontrivial.
\end{proposition}
\begin{proof}
  It suffices to find a morphism in $\Chs(K)$ whose permutation is not the identity.
  Consider the map $p : \bigvee[1,1] \to \bigvee[2]$ going along the bottom edge and then the right
  edge, and the map $q : \bigvee[1,1] \to \bigvee[2]$ going along the left edge and then the top
  edge. We have
  \[
  \begin{tikzcd}[row sep=large, column sep=large]
    \bigvee[1,1] \arrow[r, "\id"] \arrow[d, perm, "\sigma(p)"']
      & \bigvee[1,1] \arrow[d, "p"] \arrow[rd] & \\
    \bigvee[1,1] \arrow[r, "q"']
      & \bigvee[2] \arrow[r, "f"'] & K
  \end{tikzcd}
  \]
  This is a morphism of $\Chs(K)$ from $(\bigvee[1,1] \xrightarrow{f \circ p} K,\ \id)$ to
  $(\bigvee[2] \xrightarrow{f} K,\ q)$: the condition $x = \Lines(p)(y)$ is automatic because
  $\Run(\bigvee[1,1])$ is a singleton. Its permutation is
  $\sigma(p) = q_*^{-1} \circ p_* \circ \id = \{1 \mapsto 2,\ 2 \mapsto 1\}$, because $p$ traverses
  direction $1$ then direction $2$ while $q$ traverses direction $2$ then direction $1$. This has an
  inversion, so $[\sigma(p)]$ is not the identity braid.
\end{proof}

\begin{remark}[Why there is no coherent choice of run]\label{rem:no-section}
  $\pi : \Chs(K) \to \Ch(K)$ is a discrete fibration, and for a fixed $K$ it may well admit
  sections. What it does not admit is a family of sections $s_K : \Ch(K) \to \Chs(K)$ natural in
  $K$ --- and it had better not. Suppose it did. Applying naturality to the unique map $K \to Z$
  and using \Cref{cor:functorial}, $\Conc \circ s_K$ would equal $G \circ W(K)$ for the single
  functor $G := \Conc \circ s_Z$ on $\Wdg \cong \Ch(Z)$, so the braid data would be a function of
  the chain alone. For $K \subseteq \cube n$ that factors through $\Ch(\cube n)$, whose nerve is the
  contractible space $\dP(\cube n)_0^1$, and the invariant would die exactly as in
  \Cref{thm:pz-trivial}. So the content of $\Chs$ is precisely that the run \emph{cannot} be chosen
  coherently: that incoherence is the braiding. Turning this into a theorem --- fixing the class of
  $K$ over which naturality is demanded, and the sense in which the resulting invariant is trivial
  --- is \Cref{sec:next}\eqref{item:no-section}.
\end{remark}

\begin{remark}[Relation to $IB_n$]\label{rem:vs-IB}
  Two differences from the invariants of \cite[\S6]{PZ} are worth naming. Theirs is a conjugacy
  class of homomorphism out of $\pi_1(\dP(K,n), \alpha)$, so it needs a basepoint $\alpha$ and is
  only defined up to conjugacy; $\Conc$ is a functor on a category, so it needs neither. And theirs
  lands in $B_n$, the fundamental group of the \emph{unordered} configuration space, whereas
  \Cref{thm:sal} places ours over the ordered one. The $\Sigma_n$-quotient relating the two is
  \Cref{sec:next}\eqref{item:quotient}.
\end{remark}

\section{Next steps}\label{sec:next}

\begin{enumerate}
  \item\label{item:geometric} \textbf{What does $|\Chs(K)|$ compute?} \Cref{thm:ziemianski} is what
        gives $\Ch$ its meaning: its nerve is the execution space. We have no such statement for
        $\Chs$ beyond $K = \cube n$, where \Cref{thm:sal} identifies it with a configuration space.
        Is there a space of ``complexified d-paths'' on $|K|$ whose homotopy type is $|\Chs(K)|$,
        naturally in $K$? This is the main open question, and much of what follows is downstream of
        it.

  \item\label{item:quotient} \textbf{Recovering $B_n$.} $\Sigma_n$ acts on $\Sal(n)$ with quotient
        the Salvetti complex of the unordered arrangement, and the corresponding covering realises
        $P_n \le B_n$. Carrying this through on the $\Chs$ side should recover the invariants of
        \cite[\S6]{PZ} as a quotient of ours, making \Cref{rem:vs-IB} precise. The covering and the
        deck sequence are formalised in \cite{Cubical}.

  \item\label{item:no-section} \textbf{The obstruction.} Make \Cref{rem:no-section} a theorem:
        there is no section of $\pi : \Chs(-) \to \Ch(-)$ natural in $K$. The sketch given there
        reduces it to the triviality of any $\Ch$-level braid invariant on sculptures, which is
        \Cref{thm:pz-trivial}; what needs care is the quantifier over $K$ and the notion of
        triviality.

  \item\label{item:subdivision} \textbf{Subdivision and Markov moves.} The strand count of $\Conc$
        is $N = \alt(\fin_K)$, the number of events along an execution, so it changes when $K$ is
        subdivided. This is the same indexing used in \cite[\S6]{PZ}, where the invariant is defined
        per $\alpha \in \dP(K,n)_0^1$, so it is a feature of the setting rather than of this
        construction; but it does mean $\Conc$ is an invariant of the cubulation, not of the
        directed homotopy type. Looking at braid closures instead, we can relate braids with
        different numbers of strands via the Markov moves. I believe there is a connection to
        subdividing $K$, but more work is required to formalize this.

  \item\label{item:grading} \textbf{A chain complex.} The category $\Chs(K)$ is graded by the shape
        of the wedges: (the sum of the dimensions) $-$ (the length of the list). Each morphism
        should decompose into a sequence of degree $1$ morphisms. I think this is the face relation
        of a CW complex, but I'm not certain, and I also think there is a sensible boundary map
        which gives us a chain complex. Unlike the rest of the paper, none of this is formalised.

  \item\label{item:examples} \textbf{Examples.} Beyond $\cube n$ and $Z$ nothing is computed. A
        PV-program example --- the Swiss flag, say --- in which $\Conc$ distinguishes executions
        would do a lot to substantiate the claim that this is an invariant of concurrent programs.
\end{enumerate}

\printbibliography

\end{document}
